% -*- latex -*-

\label{chap:deb}
%\lstset{morecomment=[l][\color{red}]{>}
%}

\HtmlMetaDescription{This obsolete chapter describes the original
  debugging card, which has now been greatly improved and merged with
  the Debugging Front Panel subsystem. This chapter is still here
  mostly for historical reasons, as the card has been implemented and
  worked during the early days of the project.}
%\HtmlMetaGoogleDescription{}
%\HtmlMetaBanner{}
%\HtmlMetaTags{}


\begin{obsoleted}
\textbf{This peripheral has been obsoleted.} The DEB card has been merged
with the Programmer's Front Panel and updated to create the Debugging
Front Panel (DFP). This chapter is still around for historical purposes.
\end{obsoleted}

\section{Introduction}

The \gls{DEB} board originated in the early days of CFT's
\gls{Verilog} verification. A test framework would run the Verilog
simulator with various test programs, which needed to be able to print
output and make test assertions which would halt the computer when
failed.

This became a virtual device in the CFT Emulator, which allowed the
emulator to be tested on the same framework — better yet, peripherals
could also be tested, and tested much faster than in Verilog.

As a side effect, it also allowed testing of various algorithms used
in the ROM, including integer division, which also tests the
processor's overall behaviour quite well.

The success of this technique has (unsurprisingly) been great: several
serious issues with the processor and its microcode were found.

It only made sense that there should be a hardware implementation, to
bring the unit tests to the constructed CFT. As the test suite grows,
it will help verify that individual features of the processor and
computer at large work as expected.

\section{Design}

There are a number of very useful features provided by the \gls{DEB} board:

\begin{itemize}
\item It can be directed by CFT instructions to output data to the
  controlling computer.
\item It can provide information on test progress to the controlling computer.
\item It can mark failed tests to the controlling computer, and optionally halt so that the CFT's state can be inspected.
\item It can upload memory images to the computer.
\item It can examine memory locations.
\item It performs operations in I/O space, to initialise and test peripherals.
\item It can be used to reprogram the CFT's Flash ROM in-circuit.
\item It controls the run state of the machine, making it double as a
  simple virtual front panel.
\item It can be used by both humans and automated tools.
\item It provides four diagnostic LEDs (entirely user-controlled) to
  show the state of testing.
\end{itemize}

This feature set also provides a number of unforeseen benefits:

\begin{description}
\item {\bfseries Construction freedom.} Cards can now be constructed in any
  order whatsoever. The memory and IDE cards can be constructed first
  and tested to ensure they work before the processor is even
  attempted. This removes degrees of freedom and uncertainty should a
  bug be found later on. It is preferable to know that simpler devices
  (like the memory board) are working well, so issues can be tracked
  to more complex devices instead.
  
\item {\bfseries Build safety.} All the cards can be pulled out of the card
  cage, leaving \gls{DEB} board and a newly constructed peripheral
  board. The peripheral board can be tested on its own, again removing
  degrees of freedom and uncertainty. If something goes really wrong,
  the worst case scenario is that the chips on the debugging board
  need to be replaced, and the misbehaving new peripheral needs to be
  rewired slightly. This is naturally better than having to diagnose a
  fault when the set of potential failures contains over 250 ICs.

\item {\bfseries Manual operation.} With the current firmware, it is be easy
  to operate the \gls{DEB} card manually using a serial (\gls{USB})
  terminal, and get something resembling a basic programmer's front
  panel.

\item {\bfseries Flash Programming for free.} Since the \gls{DEB} board can
  write to memory and I/O devices, it can also flash the machine's
  non-volatile memory. Some flash chips (the microcode store and ALU
  tables) will have to be pulled to be programmed, but those are not
  devices that will need to be reprogrammed often.

\end{description}

\section{Theory of Operation}

The DEB board has two main parts:

\begin{description}
\item {\bfseries The CFT Bus Interface} attaches the DEB board to the CFT,
  and allows it to be addressed and used by the CFT.

\item {\bfseries The Controller Interface} interfaces the DEB board to an
  external controlling device using a TTL serial-to-USB bridge, and
  allows devices on the CFT bus (including the processor) to be
  remote-controlled, debugged and tested.
\end{description}

\subsection{Host (CFT) Interface}

\begin{figure}
\centering
\includeimage[width=\columnwidth]{figs/deb-bus-interface}\\
%\includegraphics{figs/deb-bus-interface.pdf}\\
\caption{\label{fig-deb-bus-interface}Debug board bus interface.}
\end{figure}

\begin{figure}
\centering
\includelarge[width=0.85\columnwidth]{figs/cft-deb-board-128.png}\vspace{1em}\\

\includelarge[width=0.85\columnwidth]{figs/cft-deb-board-153.png}\\
\caption{\label{fig-deb-board}A rendered image of the first draft of
  the debugging board. Note the obvious erratum: The Microchip®
  \gls{GPIO} chips are packaged in narrow 28-pin DIP, so they should
  look exactly like the \gls{MCU}.}
\end{figure}

The host interface comprises a simple address decoder and a wait-state
generator. It is shown in~\fcf{fig-deb-bus-interface}.

The address decoder decodes fully the I/O address range
\hex{03F0}–\hex{03FF} and generates a standard active-low pulse when
one of these addresses is written to.

The active low pulse asynchronously clears a D flip-flop, so its
\ns{Q} output goes high. This is inverted using an open-drain inverter
and used to drive low the \ns{WS} signal. The flip-flop is set
asynchronously by \ns{RESET} (de-asserting \ns{WS}), and synchronously
by the DEB board \gls{MCU}, which does so programmatically.

This implements the following behaviour: on reset, \ns{WS} is
guaranteed to be de-asserted. When the CFT performs an \asm{OUT \&3Fx}
instruction, the address decoder strobes the flip-flop's \ns{CLR}
line, and a wait state is asserted, freezing the computer and its buses.

Simultaneously, the strobe interrupts the \gls{MCU}, which wakes from
sleep, samples the address and data buses, acts according to the
selected operation, and strobes \ps{DEBCONT}. This sets the flip-flop,
which de-asserts \ns{WS}, and the computer resumes operation.

Please note that this wait state can be very long, perhaps in the order
or milliseconds.

\subsection{Controller Interface}

the controller interface is simple but very powerful. It comprises
three hardware parts:

\begin{description}
\item{\bfseries The \gls{MCU}} is responsible for the operation of the
  board. It is an Atmel® AVR® Atmega 168 unit with flash memory and an
  on-board SPI header for in-circuit programming.

  The \gls{MCU} controls 48 \gls{GPIO} pins, as well as the \ns{MEM},
  \ns{IO}, \ns{R} and \ns{W} conrol signals. In addition, it can stop
  and single-step the processor by driving the \ps{FPCLKEN} and
  \ps{FPUSTEP} lines.

  It also controls a bank of four LEDs which can be
    illuminated to show progress and status.

\item{\bfseries Three 16-port \gls{GPIO} chips}. Two of these are attached
  to the address and data buses (one is attached to the MSB of the
  buses, the other to the LSB of the buses). The third is for future
  expansion. A 16-pin header on the board allows signals to be
  connected to it for sampling. GPIO pins may be configured as inputs
  or outputs, reducing the chip count significantly.

  The \gls{GPIO} chips are Microchip® MCP23016 units, controlled by the
  \gls{MCU} via its I²C bus, running at 100 kHz.

  The three GPIOs are referred to by their I²C addresses or their function:
  
  \begin{itemize}
  \item \hex{20} is the address of the \gls{MSB} device, which is
    attached to \ABUSn{8}–\ABUSn{15} and \DBUSn{8}–\DBUSn{15}.
  \item \hex{21} is the address of the \gls{LSB} device, which is
    attached to \ABUSn{0}–\ABUSn{7} and \DBUSn{0}–\DBUSn{7}.
  \item \hex{23} is the address of the user device, which can be used to
    sample or control arbitrary lines.
  \end{itemize}

\item{\bfseries The serial-to-USB bridge} connects the \gls{MCU} to a
  testing controller via \gls{USB}. This allows high communication
  speeds — by default, the DEB card communicates at 921,600 bps, while
  ensuring near-universal compatibility with computing devices.
\end{description}

The \gls{MCU} firmware controls the hardware assets programmatically,
so behaviours can be changed, and features may be upgraded as
required.

\section{Operation}
\label{sec-deb-firmware}

When powered on or reset, the board goes through a brief
self-diagnostic phase, ensuring that the \gls{GPIO} chips are in good
working order. The LED testing lights should illuminate, then turn off
in sequence, until none remain lit. This process should last
no longer than around three seconds.

If the board stops before completing initialisation,
\tcf{tab-deb-diag} can be consulted for the meaning of the diagnostic
LEDs.

If all devices pass power-on tests, all the LEDs should be off.

If a \gls{GPIO} is reported failed, the fault may be transient. The
MCP23016 \gls{GPIO} has no reset signal, and its error recovery is
less than perfect — if an I²C or protocol error is encountered, it can
hang. Resetting the \gls{DEB} board is not enough to clear the error
condition. Power cycling is recommended.

Diagnostics are also indicated via the serial connection. Upon
power-up or reset, the following should appear:

\begin{debcode}
CFT DEB Board Booting.
GPIO diag: 20 21 23 Ready.
201 Version: 0.5
201 BufSize: 104
202 (c) 2012 Alexios Chouchoulas <alexios@bedroomlan.org>
202 Licensed under the GNU Public License v.2.
202 http://www.bedroomlan.org/cft
[running]>
\end{debcode}

\newcommand\led[1]{%
  \ifxetex
    \begin{tikzpicture}
      \draw[line width=1pt, fill=#1] (0,0) circle(4pt);
    \end{tikzpicture}
  \else
    \HCode{<span class="led #1">&\#9679;</span>}
  \fi
}

\begin{table}
  \zebra
  \centering
  \begin{tabular}{ccccl}
    \noalign{\smallskip}\hline\noalign{\smallskip}
    A & B & C & D & Meaning \\
    \noalign{\smallskip}\hline\noalign{\smallskip}
    \led{r1} & \led{g1} & \led{g1} & \led{g1} & Early hardware failure. \\
    \led{r1} & \led{g1} & \led{g1} & \led{g0} & MSB \gls{GPIO} fault. \\
    \led{r1} & \led{g1} & \led{g0} & \led{g0} & LSB \gls{GPIO} fault. \\
    \led{r1} & \led{g0} & \led{g0} & \led{g0} & User \gls{GPIO} fault. \\
    \led{r0} & \led{g0} & \led{g0} & \led{g0} & Device is ready. \\
    \noalign{\smallskip}\hline\noalign{\smallskip}
  \end{tabular}
\caption{\label{tab-deb-diag} Meaning of the \gls{DEB} diagnostic
  lights, should the board stop during initialisation.}
\end{table}


\section{Interactive Operation — The Virtual Front Panel}

The firmware presents an interactive user interface intended for both
humans and machines, in the style of the SMTP, FTP and HTTP
protocols. The user enters short commands, optionally followed by one
or more arguments, and the board responds with a three digit response
code and a human-readable message. A human ignores the response code;
a machine ignores the textual message.

A rudimentary line editor is provided for command lines:

\begin{description}
\item{\bfseries Backspace or DEL} keys delete the last character in the
  current command.
\item{\bfseries Ctrl-C or Ctrl-X} keys abort the current command, and may
  also be used to abort a long-running process.
\item{\bfseries Newline or Line-feed} keys submit the current command for
  processing.
\end{description}

Help may be obtained by using the \cftin{?} or \cftin{help} commands,
which print out something like this:

\begin{debcode}
[running]> help

201 Available commands:
201  ? -- Show help.
201  help -- Show help.
201  ver? -- Show version.
201  buf? -- Show size of command buffer.
201  c? -- Show current word count.
201  ob -- Set binary for dump.
201  ot -- Set text for dump.
201  leds NIBBLE -- Set diagnostic LEDs (0-F).
201  halt -- Halt the CFT & become bus master.
201  run -- Resume the CFT.
201  rst! -- Reset the CFT.
201  e0 -- Local echo off.
201  e1 -- Local echo on.
201  hof0 -- Halt on FAIL.
201  hof1 -- Continue on FAIL.
201  a WORD -- Set address.
201  c WORD -- Set count.
201  r -- Read word from memory.
201  w WORD ... -- Write words to memory, step address.
201  in WORD -- Read from I/O address WORD.
201  out ADDR WORD -- Write WORD to I/O ADDR.
201  d -- Dump count words starting at addr, step addr.
201  clr! -- Reset the DEB board.
201 Consult documentation for details.
\end{debcode}

The \gls{DEB} board has two modes of operation:

\begin{description}

\item{\bfseries Run mode}. In this mode, the CFT is executing its program as
  per normal. It is not possible to modify the computer's state in
  this mode, except to stop it or reset it. The \cftin{run} command
  (described below) activates this mode.

\item{\bfseries Stop mode}. In this mode, the \gls{DEB} board has master
  access to the computer's buses and can inspect and modify memory and
  I/O locations. The \cftin{halt} command (described below) activates
  this mode and unlocks the full \gls{DEB} command set.

\end{description}

The following is a brief description of available \gls{DEB} commands.

\begin{description}
\item{\bfseries\cftin{?} or \cftin{help}} Lists available commands and their
  brief descriptions.

\item{\bfseries\cftin{ver?}} Prints out the current firmware version. The ‘+t’
  signifies that text dumping is available, which is always the case
  on the Atmega 168-based board.\footnote{This originated in the
    author's Flash Programmer project, which was coded for both
    Atmega 168 and Atmega 8. The latter has 8 kBytes of flash
    memory, and text-based dumping would not fit.}

\begin{debcode}
[running]> ver?
201 Version: 0.5+t
\end{debcode}

\item{\bfseries\cftin{buf?}} Displays the size of the input line buffer {\bfseries
  in hexadecimal}. This affects the maximum number of values that can
  be written to memory on one line.

\begin{debcode}
[running]> buf?
201 BufSize: 200
\end{debcode}

\item{\bfseries\cftin{c?}} Displays the current word count {\bfseries in
  hexadecimal}. This is the size of the block displayed by the \cftin{d}
  (dump) command. The default count is \hex{80}, i.e. 128 words.

\begin{debcode}
[running]> c?
201 Count: 0080
\end{debcode}

\item{\bfseries\cftin{ob}} Makes the \cftin{d}
  (dump) command dump in binary. The default is text dumping.

\begin{debcode}
[running]> ob
206 Output: bin
\end{debcode}

\item{\bfseries\cftin{ot}} Makes the \cftin{d}
  (dump) command dump in text, which is also the default.

\begin{debcode}
[running]> ot
205 Output: text
\end{debcode}

\item{\bfseries\cftin{leds} NYBBLE} Sets the bit pattern of the diagnostic
  LEDs. \gls{NYBBLE} is a hexadecimal value in the range
  \hex{0}–\hex{F} (four bits). Bits 0–3 represent LEDs A–D in order. LED D is the
  red LED. A bit value of \bin{1} turns on a LED; a value of \bin{0}
  turns it off. To turn on all LEDs but LED A (\bin{1110} = \hex{E}):

\begin{debcode}
[running]> leds e
209 LEDs: e
\end{debcode}

\item{\bfseries\cftin{halt}} Stops the computer and unlocks the full \gls{DEB}
  command set. The mode changes to halted, and the prompt changes to
  show the current address:

\begin{debcode}
[running]> halt
306 Host halted
[halted] 0000> halt
505 Already halted.
[halted] 0000> 
\end{debcode}


\item{\bfseries\cftin{run}} Locks the \gls{DEB} command set and starts the
  computer.  The mode changes to running, and the prompt changes to
  reflect this:

\begin{debcode}
[halted] 0000> run
305 Host running
[running]> run
506 Already running.
[running]> 
\end{debcode}


\item{\bfseries\cftin{rst!}} Resets the CFT processor. The \gls{DEB} board is
  not affected by this. No confirmation questions are asked.

\begin{debcode}
[halted] 0000> rst!
307 Host reset
\end{debcode}


\item{\bfseries\cftin{e0}} Disables local echo. This is useful for certain
  remote terminals that have custom line editors and/or their own
  echo, which make typed characters appear twice. The default is to
  perform local echo (\cftin{e1}).

\begin{debcode}
[halted] 0000> e0
208 Echo: off
\end{debcode}


\item{\bfseries\cftin{e1}} Enables local echo. Use this command if typed
  characters are invisible. This is the default behaviour.

\begin{debcode}
[halted] 0000> e1
207 Echo: off
\end{debcode}


\item{\bfseries\cftin{hof0}} Disables halt-on-\asm{FAIL} mode. When the CFT
  issues the \asm{FAIL} instruction, the computer will be allowed to
  keep running. This is the default behaviour.

\begin{debcode}
[halted] 0000> hof0
211 Halt on FAIL: off
\end{debcode}


\item{\bfseries\cftin{hof1}} Enables halt-on-\asm{FAIL} mode. When the CFT
  issues the \asm{FAIL} instruction, the computer will be halted (and
  the full \gls{DEB} command set unlocked) to allow state
  inspection. The default behaviour is \cftin{hof0}.

\begin{debcode}
[halted] 0000> hof1
210 Halt on FAIL: on
\end{debcode}


\item{\bfseries\cftin{a} WORD} Sets the current address to WORD, which is a
  16-bit value expressed {\bfseries in hexadecimal}. The current address
  affects memory reads (using the \cftin{d} command) and memory writes
  (using the \cftin{w} command). If the \gls{DEB} board is in the
  halted mode, the prompt also indicates the current address:

\begin{debcode}
[halted] 0000> a 1234
203 Address: 1234
[halted] 1234> 
\end{debcode}


\item{\bfseries\cftin{c} COUNT} Sets the current count to COUNT, which is a
  16-bit value expressed {\bfseries in hexadecimal}. Please note that this
  makes it impossible to, for instance, dump 65,536 words (the full
  CFT logical address space) because the largest acceptable value for
  COUNT is \hex{FFFF} or 65,535.

\begin{debcode}
[halted] 0000> c 8000
204 Count: 8000
\end{debcode}


\item{\bfseries\cftin{r}} Reads one word from memory space at the current
  address. The read word is reported {\bfseries in hexadecimal}. The current
  address is incremented by one. This command only works with the
  \gls{DEB} board in halted mode.

\begin{debcode}
[running]> a 1234
203 Address: 1234
[running]> r
503 Halt host first.
[running]> halt
306 Host halted
[halted] 1234> r
304 Value: DEAD
[halted] 1235> r
304 Value: BEEF
[halted] 1236> r
304 Value: 63FF
[halted] 1237> 
\end{debcode}


\item{\bfseries\cftin{w} WORD …} Writes one or more words to space starting
  with the current address. Words are 16-bit values expressed {\bfseries in
    hexadecimal}. The current address is incremented after each
  write. This command only works with the \gls{DEB} board in halted
  mode. A simple checksum (a 32-bit sum of all input values expressed
  {\bfseries in hexadecimal} and padded to 8 digits) is reported.

\begin{debcode}
[running]> a 1234
203 Address: 1234
[running]> w dead beef 63ff
503 Halt host first.
[running]> halt
306 Host halted
[halted] 1234> w dead beef 63ff
301 Done
303 Checksum: 0002019b
[halted] 1237> 
\end{debcode}


\item{\bfseries\cftin{in} WORD} Reads one 16-bit value from I/O space address
  WORD and reports it in hexadecimal, padded to 4 digits. This command
  only works with the \gls{DEB} board in halted mode.

\begin{debcode}
[running]> in 23
503 Halt host first.
[running]> halt
306 Host halted
[halted] 0000> in 23
304 Value: 00FF
[halted] 0000> 
\end{debcode}


\item{\bfseries\cftin{out} ADDR WORD} Writes the 16-bit WORD to I/O space
  address ADDR. This command only works with the \gls{DEB} board in
  halted mode. The current address is not modified.

\begin{debcode}
[running]> out 16f 1234
503 Halt host first.
[running]> halt
306 Host halted
[halted] 0000> out 16f 1234
301 done
[halted] 0000> 
\end{debcode}


\item{\bfseries\cftin{d}} Dumps a block of COUNT (see \cftin{c} and \cftin{c?}
  commands above) words starting at the current address. The current
  address is incremented by COUNT. This command only works with the
  \gls{DEB} board in halted mode. In text mode (see \cftin{ot}), the
  command works like this:

\begin{debcode}
[running]> c 20
204 Count: 0020
[running]> d
203 Address: 0000
204 Count: 0020
300 Dumping
503 Halt host first.
[running]> halt
306 Host halted
[halted] 0000> d
203 Address: 0000
204 Count: 0020
300 Dumping
0000 | 1234 1234 1234 1234 1234 1234 1234 1234 | 44444444 | 000091a0
0008 | 1234 1234 1234 1234 1234 1234 1234 1234 | 44444444 | 00012340
0010 | 1234 1234 1234 1234 1234 1234 1234 1234 | 44444444 | 0001b4e0
0018 | 1234 1234 1234 1234 1234 1234 1234 1234 | 44444444 | 00024680
301 Done
303 Checksum: 00024680
[halted] 0020>
\end{debcode}

The dump is in four columns. The leftmost column is the current
address. The second column lists eight 16-bit hexadecimal word values
(padded to four digits), starting at the current address. The third
column shows ASCII characters by examining the least significant seven
bits of the value. Unprintable characters are replaced with full stops
(\cftin{.}). The rightmost column shows a running checksum of the values
listed so far (including previous lines).

At the end of the dump, the final checksum is reported.

In binary mode, the newline following the ‘\cftin{300 Dumping}’ message
signals the start of the dump. 2×COUNT 8-bit bytes are
transmitted. Each pair of bytes is one value from memory. The
\gls{MSB} is transmitted first. The receiver should read exactly
2×COUNT bytes and revert to parsing the ‘\cftin{301 Done}’ result. As
with the text dump, the block checksum is reported.

\begin{debcode}
[halted] 0000> c 20
204 Count: 0020
[halted] 0000> d
203 Address: 0000
204 Count: 0020
300 Dumping
(64 bytes of binary data follow)
301 Done
303 Checksum: 00024680
[halted] 0020>
\end{debcode}



\item{\bfseries\cftin{clr!}} Performs a warm reset of the \gls{DEB} board,
  reverting all settings to their defaults. This does not affect
  operation of the CFT computer itself, nor does it change the current
  mode run/stop mode. Thus, a halted computer will remain halted after
  the warm reset.

\begin{debcode}
[halted] 6120> clr!
GPIO diag: 20 21 23 Ready.
201 Version: 0.5+t
201 BufSize: 200
202 (c) 2012 Alexios Chouchoulas <alexios@bedroomlan.org>
202 Licensed under the GNU Public License v.2.
202 http://www.bedroomlan.org/cft
[halted] 0000>
\end{debcode}

\end{description}

\subsection{Result Codes}

The \gls{DEB} firmware reports three-digit result codes in
decimal. These are broken down into four categories:

\begin{description}
\item{\bfseries 200–299} Informational messages and acknowledgements. These
  do not affect the operation of the computer.
  \begin{description}
  \item{\bfseries\cftout{201}} Supplemental information, boot-up banners.
  \item{\bfseries\cftout{202}} Copyright and licensing notices.
  \item{\bfseries\cftout{203 Address:}} The current address is the last field or last four
    characters of the line.
  \item{\bfseries\cftout{204 Count:}} The word count the last field or last four characters
    of the line.
  \item{\bfseries\cftout{205 Output: text}} Dump output mode has been switched to text.
  \item{\bfseries\cftout{206 Output: bin}} Dump output mode has been switched to binary.
  \item{\bfseries\cftout{207 Echo: on}} Local echo activated.
  \item{\bfseries\cftout{208 Echo: off}} Local echo deactivated.
  \item{\bfseries\cftout{209 LEDS: on}} LED output pattern set. The pattern is
    the last space-delimited field (one character) on the line.
  \item{\bfseries\cftout{210 Halt on FAIL: on}} Halt-on-FAIL enabled.
  \item{\bfseries\cftout{211 Halt on FAIL: off}} Halt-on-FAIL disabled.
  \end{description}
  
\item{\bfseries 300–399} Messages pertaining to processes that access or
  modify computer state.
  \begin{description}
  \item{\bfseries\cftout{300 Dumping}} A memory dump follows.
  \item{\bfseries\cftout{301 Done}} Process is finished.
  \item{\bfseries\cftout{302 Aborted}} A process was aborted by the user.
  \item{\bfseries\cftout{303 Checksum: }} The last space-delimited field on
    the line (eight characters) is a 32-bit hexadecimal representation of the
    checksum of the last block of data sent or received.
  \item{\bfseries\cftout{304 Value:}} The last space-delimited field on the
    line (four characters) is a 16-bit hexadecimal representation of
    the value read from memory, as requested.
  \item{\bfseries\cftout{305 Host running}} The CFT computer was started or
    restarted. \gls{DEB} commands that use the computer buses are now
    locked out.
  \item{\bfseries\cftout{306 Host halted}} The CFT computer was halted. The
    full \gls{DEB} command set is now unlocked.
  \item{\bfseries\cftout{307 Host reset}} The CFT computer was reset.
  \end{description}
  
\item{\bfseries 400–499} Warnings.
  \begin{description}
  \item{\bfseries\cftout{401 Warning: write will wrap around}} If COUNT
    words are written at the current address, the address will wrap
    around to \hex{0000} during the block write.
  \end{description}
  
\item{\bfseries 500–599} Errors.
  \begin{description}
  \item{\bfseries\cftout{500 Unknown command}} The last command was not recognised.
  \item{\bfseries\cftout{501 Bad value}} An argument was not a valid hexadecimal number.
  \item{\bfseries\cftout{502 Count must be multiple of 8}} For proper
    behaviour of dumping, the word count should have its least
    significant three bits clear, i.e. be a multiple of eight.
  \item{\bfseries\cftout{503 Halt host first}} A locked-out command was
    issued. The computer must be halted for these commands to become
    available.
  \item{\bfseries\cftout{504 Bus chatter}} For safety, CFT buses are sampled
    before driving them. If changing values are seen, this message
    appears and the command is aborted to avoid bus contention. This
    is an indication of something wrong with the bus connections.
  \item{\bfseries\cftout{505 Already halted}} The computer is already
    halted. The \cftin{halt} command had no effect.
  \item{\bfseries\cftout{506 Already running}} The computer is already
    running. The \cftin{run} command had no effect.
  \end{description}
\end{description}

The following is a list of all the result codes in use by the firmware.

\section{Automated Operation — The DEB Tool}

The \gls{DEB} board may be used remotely by a controlling computer, to
automate unit testing of the CFT. The tool allows memory images to be
downloaded to the DEB card and written to CFT memory, then executed,
and their output captured for analysis by the controlling computer.

\todo{The DEB controller-side tool has yet to be written.}

\section{Operation from the CFT}

The mapping of the \gls{DEB} board to the I/O space provides sixteen
extended CFT instructions with unusually long names due to their
unique history. The following is a list of these I/O addresses:

\begin{extcmd}{PRINTA}{OUT R \&3F0}{57F0}{c-wvehf}{3F0}{DEB, PRINT Address}             
  The current value of \A{} is output from the
  \gls{DEB} board to the controlling computer in the format
  \cftout{XXXX (label+N)}, where \cftout{XXXX} is the hexadecimal
  value of \A{} padded to four digits, \cftout{label} is the most
  recent \gls{Assembly} source label, and \cftout{n} is the offset (in
  words) from that label. If a label is not available, ‘\cftout{XXXX
    (?)}’ is output instead. Labels are only available on the emulator
  with a map file loaded.%
\end{extcmd}

\begin{extcmd}{PRINTC}{OUT R \&3F1}{57F1}{c-wvehf}{3F1}{DEB, PRINT Character}
  The \gls{LSB} of the current value of \A{} is output from the
  \gls{DEB} board to the controlling computer as an 8-bit
  character.
\end{extcmd}

\begin{extcmd}{PRINTD}{OUT R \&3F2}{57F2}{c-wvehf}{3F2}{DEB, PRINT Decimal}
  The current value of \A{} is output from the \gls{DEB} board to the
  controlling computer as a decimal signed, \gls{Twos Complement}
  integer in the range -32,768–32,767. There is no zero-padding.
\end{extcmd}

\begin{extcmd}{PRINTU}{OUT R \&3F3}{57F3}{c-wvehf}{3F3}{DEB, PRINT Unsigned integer}
  The current value of \A{} is output from the \gls{DEB} board to the
  controlling computer as a decimal unsigned integer in the range
  0–65,535. There is no zero-padding.
\end{extcmd}

\begin{extcmd}{PRINTH}{OUT R \&3F4}{57F4}{c-wvehf}{3F4}{DEB, PRINT Hexadecimal}
  The current value of \A{} is output from the \gls{DEB} board to the
  controlling computer as a hexadecimal unsigned integer in the range
  \hex{0000}–\hex{FFFF}. The number is zero-padded to four digits. The
  case of the alphabetic digits is purposefully unspecified. Parsers
  should understand both lower and upper case.
\end{extcmd}

\begin{extcmd}{PRINTB}{OUT R \&3F5}{57F5}{c-wvehf}{3F5}{DEB, PRINT Binary}
  The current value of \A{} is output from the \gls{DEB} board to the
  controlling computer in its raw binary representation, in the range
  \hex{0000000000000000}–\hex{1111111111111111}. The number is
  zero-padded to sixteen bits.
\end{extcmd}

\begin{extcmd}{PRINTSP}{OUT R \&3F6}{57F6}{c-wvehf}{3F6}{DEB, PRINT SPace}
  A single space (ASCII codepoint 32) is output from the \gls{DEB} board
  to the controlling computer. This simplifies debugging and testing,
  since none of the \gls{DEB} instructions space-pad their output.
\end{extcmd}

\begin{extcmd}{PRINTNL}{OUT R \&3F7}{57F7}{c-wvehf}{3F7}{DEB, PRINT New Line}
  A single new line (ASCII codepoint 10) is output from the \gls{DEB} board
  to the controlling computer.
\end{extcmd}

\begin{extcmd}{DEBUGON}{OUT R \&3F8}{57F8}{c-w-e-f}{3F8}{DEB, Enable Assembly tracing}
  Only available in the emulator. Enables symbolic \gls{Assembly} tracing of all
  instructions executed.
\end{extcmd}

\begin{extcmd}{DEBUGOFF}{OUT R \&3F9}{57F9}{c-w-e-f}{3F9}{DEB, Disable Assembly tracing}
  Only available in the emulator. Disables symbolic \gls{Assembly} tracing of all
  instructions executed.
\end{extcmd}

\begin{extcmd}{DUMP}{OUT R \&3FA}{57FA}{c-w-e-f}{3FA}{DEB, DUMP State}
  Not currently available in the hardware implementation. Dumps the
  current state of the computer in an implementation-specific way
  intended for post-mortem debugging.
\end{extcmd}

\begin{extcmd}{PRINTHI}{OUT R \&3FB}{57FB}{c-w-ehf}{3FB}{DEB, PRINT HIgh Word}
  Stores the current value of the \A{} register in a temporary register
  on board the \gls{DEB} board. Meant to be precede the
  \asm{PRINTLO} instruction to print out 32-bit hexadecimal values.
\end{extcmd}

\begin{extcmd}{PRINTLO}{OUT R \&3FC}{57FC}{c-w-ehf}{3FC}{DEB, PRINT LOw Word}
  Prints out the 32-bit integer made by combining the last value stored
  by \asm{PRINTHI}, shifted 16 bits to the left, and the current value
  of \A{}. The number reported is in the range
  \hex{00000000}–\hex{FFFFFFFF} and is zero padded to eight digits. The
  case of the alphabetic digits is purposefully unspecified. Parsers
  should understand both lower and upper case.

  The initial value of the \asm{PRINTHI} register is undefined — care
  should be taken to always issue a \asm{PRINTHI} before a
  \asm{PRINTLO}.
\end{extcmd}

\begin{extcmd}{SUCCESS}{OUT R \&3FE}{57FE}{c-wvehf}{3FE}{DEB, Assert SUCCESSful test checkpoint}
  Registers a successful test checkpoint with the test harness. These
  may be used to provide progress reporting as well as to ensure an
  exact, certain number of checkpoints is reached by the code.

  To mark the checkpoint, the sequence ‘\cftout{[ok]}’ is transmitted to
  the controlling computer.
\end{extcmd}

\begin{extcmd}{FAIL}{OUT R \&3FF}{57FF}{c-wvehf}{3FF}{DEB, FAILed test assertion}
  Registers a failed test assertion with the test harness. This
  is used to provide information to the controlling computer.

  The failure is denoted by sending the sequence ‘\cftout{[fail]}’ to the
  controlling computer.

  On the \gls{DEB} board hardware implementation, \asm{FAIL} illuminates
  the red LED (LED D). If Halt-on-FAIL mode is enabled, the processor is halted.

  On the Verilog simulator and CFT Emulator, \asm{FAIL} always
  terminates the Verilog or emulator session, and usually dumps the full
  state of the CFT and its peripherals for post-mortem analysis
  (identical to the side-effects of the \asm{DUMP} instruction).
\end{extcmd}


\section{Notes}

\notes{../../notes/debug-board.txt}


\section{Schematics}

The following page shows the full schematic drawing of the \gls{DEB} board.

\includeschns{45}{DEB Board}{sch:deb}
